In this section and the next we consider the firm's production in the short run.\footnote{
All proofs are collected in Appendix~\ref{app:omitted_proofs}. For the reader's convenience, Table~\ref{tab:notation} in Appendix~\ref{app:theory_tables} summarizes the notation used throughout the theory sections.
}
A firm's production consists of a sequence of steps $\mathcal{S}=(s_1,\ldots,s_m)$, every one of which must be completed to produce a unit of output.\footnote{
We hold the set of production steps fixed and abstract from the creation or disappearance of steps in our model.
}
Completing a step takes worker time and demands skill, and both depend on the mode in which the step is executed.
The firm deploys AI by assigning each step to a human or to AI, and that assignment determines which tasks the worker actually handles.

In the short run, the set of steps assigned to each worker is fixed (i.e., job boundaries are already drawn); workers have acquired the skills their roles require; and the firm has hired them at wages agreed in advance.
The firm's only margin of adjustment is AI deployment, that is, which steps it assigns to AI.
That choice sets how much time each worker spends completing steps, and hence the wage bill the firm pays.
Given these short-run assumptions, we assume without loss of generality that a single worker holds the entire production sequence, and we normalize the wage to $1$.
Section~\ref{sec:longrun} relaxes these, letting the firm divide the steps among workers so that the skill a job demands, and hence the wage it commands, becomes endogenous.

\begin{figure}[!t]
  \centering
  \caption{Bundling Steps into Tasks via AI Strategy}
  \label{fig:steps_tasks}
  % Steps-to-Tasks Visualization.
% This is the top two layers of plots/TikZ_visualization/hierarchical_production_illustration.tex,
% cut at the task layer and kept identical to it in scale, node distance, box
% styles, and colors, so that a step, a task, and an aggregation arrow read the
% same in both figures. The only omissions are the job layer and the red snaked
% hand-off arrows, neither of which exists before jobs are introduced.
\begin{adjustbox}{center,trim=0pt 0pt 0pt 0pt,clip,margin=-2cm 0cm}
  \begin{tikzpicture}[scale=0.73, transform shape,
    node distance=0.35cm and 0.91cm,
    every node/.style={rectangle, rounded corners, draw, align=center, minimum width=2cm, minimum height=1cm},
    manual/.style={fill=gray!10},
    automated/.style={fill=green!30},
    augmented/.style={fill=orange!30},
    dashedbox/.style={draw, rectangle, rounded corners, dashed, inner sep=0.25cm, line width=1.25pt},
    humanbox/.style={rectangle, rounded corners, draw, align=center, minimum width=2.25cm, minimum height=1cm, fill=olive!20},
    aibox/.style={rectangle, rounded corners, draw, align=center, minimum width=2.25cm, minimum height=1cm, fill=blue!20},
    >=latex
  ]

    % Top layer (Steps)
    \node[manual] (S1) {Step 1\\(Manual)};
    \node[augmented, right=of S1] (S2) {Step 2\\(Augmented)};
    \node[manual, right=of S2] (S3) {Step 3\\(Manual)};
    \node[automated, right=of S3] (S4) {Step 4\\(Automated)};
    \node[automated, right=of S4] (S5) {Step 5\\(Automated)};
    \node[augmented, right=of S5] (S6) {Step 6\\(Augmented)};
    \node[manual, right=of S6] (S7) {Step 7\\(Manual)};

    % Arrows between steps
    \draw[->, thick] (S1) -- (S2);
    \draw[->, thick] (S2) -- (S3);
    \draw[->, thick] (S3) -- (S4);
    \draw[->, thick] (S4) -- (S5);
    \draw[->, thick] (S5) -- (S6);
    \draw[->, thick] (S6) -- (S7);

    % Dashed boxes grouping consecutive steps into tasks
    \node[dashedbox, draw=olive, fit=(S1)] (DS1) {};
    \node[dashedbox, draw=blue, fit=(S2)] (DS2) {};
    \node[dashedbox, draw=olive, fit=(S3)] (DS3) {};
    \node[dashedbox, draw=blue, fit=(S4)(S5)(S6)] (DS4-6) {};
    \node[dashedbox, draw=olive, fit=(S7)] (DS7) {};

    \path (S1.west) -- (S7.east) coordinate[midway] (CenterTop);

    % Bottom layer (Tasks)
    \node[humanbox, below=2.4cm of CenterTop, xshift=-6cm] (T1) {Task 1\\(Human)};
    \node[aibox, right=of T1] (T2) {Task 2\\(AI Chain)};
    \node[humanbox, right=of T2] (T3) {Task 3\\(Human)};
    \node[aibox, right=of T3] (T4) {Task 4\\(AI Chain)};
    \node[humanbox, right=of T4] (T5) {Task 5\\(Human)};

    % Arrows between tasks
    \draw[->, thick] (T1) -- (T2);
    \draw[->, thick] (T2) -- (T3);
    \draw[->, thick] (T3) -- (T4);
    \draw[->, thick] (T4) -- (T5);

    % Aggregation arrows, colored to match the dashed box each task comes from
    \draw[->, thick, olive] (DS1.south) -- (T1.north);
    \draw[->, thick, blue]  (DS2.south) -- (T2.north);
    \draw[->, thick, olive] (DS3.south) -- (T3.north);
    \draw[->, thick, blue]  (DS4-6.south) -- (T4.north);
    \draw[->, thick, olive] (DS7.south) -- (T5.north);

  \end{tikzpicture}
\end{adjustbox}

  \par\vspace{\fignotesskip}
\begin{minipage}{1\linewidth}
\begin{spacing}{0.2}
{\fontsize{9.5pt}{10pt}\selectfont Notes: The top layer shows seven production steps, each completed in manual (gray), automated (green), or augmented (orange) mode; dashed boxes group contiguous steps into tasks. The bottom layer shows the resulting tasks: a run of automated steps ending in an augmented step forms an AI chain (blue; here steps 4--6 become Task~4), and a lone augmented step is an AI chain of length one (here step~2 becomes Task~2), while a step performed manually on its own is a human task (olive). Horizontal arrows show the production flow within each layer, and each colored vertical arrow shows which dashed block of steps becomes which task.}
\end{spacing}
\end{minipage}
\end{figure}



\subsection{Steps}
\label{sec:model.steps}

Steps are the smallest unit of work in our framework, and each is executed in one of three modes: \emph{manually}, \emph{augmented}, or \emph{automated}.
If the firm assigns step $i$ to manual execution (denoted by \manualLetter), a human worker spends time \manualTime{i} completing the step without AI.
Alternatively, the firm may assign the step to AI (denoted by \AIletter), in which case a worker spends time \AItime{i} prompting and verifying the output of one AI attempt at it.\footnote{
This could, for example, represent the time spent specifying what the step requires and then checking whether what the AI returns meets it.
}
Such an attempt succeeds independently with some (step-dependent) probability $q_i=\alpha^{d_i}$, where $\alpha$ is the quality of the general-purpose AI technology and $d_i$ represents how ``difficult'' step $i$ is for AI, and it must be repeated until it succeeds. 
Since the worker repeats prompting and verifying AI's output until success, the expected time cost of verification across trials is $\AItime{i}/q_i$.
Definitions~\ref{def:manual_step} and \ref{def:augmented_step} formalize these two modes of step execution.

\begin{definition}[Manual Step]
\label{def:manual_step}
A step $i$ is performed manually if it is executed entirely by a human worker without AI assistance.
The time cost of a manual step is denoted by \manualTime{i}.
\end{definition}

\begin{definition}[Augmented Step]
\label{def:augmented_step}
A step $i$ is augmented if it is executed by the AI, after which a human worker verifies its output.
The expected time cost of verifying a single augmented step is $\AItime{i}/q_i$, where $q_i = \alpha^{d_i} \in (0,1]$ is the AI's probability of successfully completing the step, increasing in the quality $\alpha$ of the general-purpose AI technology and decreasing in the step's difficulty $d_i$ for AI.
\end{definition}

In the top layer of Figure~\ref{fig:steps_tasks}, Step~1, for example, is executed manually and hence is a manual step, whereas Step~2 is an instance of an augmented step performed by the AI and verified by a human.

There is a third mode of executing a step, namely automation, that makes AI more than a faster way of doing the same work, but it is not freely available.
Because no human handles an automated step, nothing checks its output, so that output must feed directly into another AI-executed step rather than into a worker.
Automation therefore arises only inside a contiguous run of AI-executed steps whose last step is augmented, which is where the human re-enters and verifies what the run produced.
We call such a run an \emph{AI chain}.

Formally, consider an AI chain spanning steps $(s_\ell, \dotsc, s_r)$.
The worker verifies only the output of its final step $s_r$, since all earlier steps run automatically ``under the hood.''
Verifying one attempt at the chain therefore costs the same as verifying step $s_r$ on its own, namely time \AItime{r} per attempt.
An attempt succeeds only if AI completes every step of the chain, which, assuming independent failures, occurs with probability $\prod_{i=\ell}^{r} q_i$; the chain is retried until it succeeds, for a total expected time of $\AItime{r}/\prod_{i=\ell}^{r} q_i$.
Since $q_i = \alpha^{d_i}$, the chain succeeds with probability $\alpha^{\sum_{i=\ell}^{r} d_i}$, so $\sum_{i=\ell}^{r} d_i$ is the chain's total difficulty, aggregating additively over its constituent steps.
A chain may consist of a single step, with $\ell = r$: having no automated predecessors, it is simply an augmented step, so an augmented step is an AI chain of length one, a convention we adopt throughout the paper.

Definitions~\ref{def:automated_step} and \ref{def:ai_chain} formalize these two additional concepts.

\begin{definition}[Automated Step]
\label{def:automated_step}
A step is automated if it is executed by AI end-to-end, with no human verifying its output, so that its direct human time cost is zero.%
\footnote{That execution of automated steps does not incur any cost to the firm is consistent with the AI subscription business model, where the firm pays an upfront fixed cost to an AI provider and uses AI tools at zero marginal cost afterward. The main forces discussed in the model remain at play even if we assume that each time the firm deploys AI on a step it has to incur a cost. We deliberately avoid modeling this to prevent complicating the model.}
\end{definition}

\begin{definition}[AI Chain]
\label{def:ai_chain}
An AI chain is a contiguous block of one or more sequential steps executed by AI, in which all steps but the final one are automated and the final step is augmented.
An AI chain spanning steps $(s_\ell,\dots,s_r)$ has steps $(s_\ell,\dots,s_{r-1})$ automated and step $s_r$ augmented, with expected time cost $\AItime{r}/\prod_{i=\ell}^{r} q_i$.
\end{definition}
In Figure~\ref{fig:steps_tasks}, Steps~4--6 form an AI chain. The human prompts the AI for, and verifies the output of, only the chain's augmented step, which for Steps~4--6 is Step~6; Steps~4 and~5 are performed without any direct human intervention and are hence automated.
Step~2 forms an AI chain (of length one) on its own as well, containing only a single augmented step.
Next, we introduce how the firm's placement of these chains determines the tasks its workers hold.



\subsection{Tasks}
\label{sec:model.tasks}

If steps are the smallest unit of work, tasks are the smallest unit of \emph{human} work.
The distinction follows from the execution modes of steps above. 
A worker is involved in a step only when they perform it manually or verify what the AI returns at the end of a chain, whereas an automated step runs inside the chain entirely beneath their awareness.
Every point at which a worker executes or verifies work therefore opens a task for them, and whatever the AI has run automatically since the last such point is absorbed into it.
The steps of a workflow map into tasks accordingly: a manual step is a task on its own, and an entire AI chain is a single \emph{composite} task, delegated to the AI at its start and received back at its augmented endpoint.
However many steps a chain spans, it costs the worker the same (per attempt) to verify, which is the sense in which chaining changes not just how fast work is done but how much of it a human ever handles.

The firm's choice of where to place chains is therefore what creates the tasks its workers hold, and we record that choice as an \emph{AI deployment strategy}.

\begin{definition}[AI Deployment Strategy]
\label{def:ai_strategy}
An AI deployment strategy (or AI strategy for short) is a partition of the step sequence $\mathcal{S}=(s_1,\ldots,s_m)$ into contiguous blocks, each of which is either a single manual step or an AI chain.
The blocks of the partition are the firm's tasks, forming the task sequence $\T = (T_1, \dotsc, T_n)$, and each task $T_b$ carries the expected time cost \timecost{b} of the mode that executes it: \manualTime{i} for a manual step $s_i$, and $\AItime{r}/\prod_{i=\ell}^{r} q_i$ for a chain spanning $(s_\ell, \dotsc, s_r)$.
\end{definition}

The AI strategy in Figure~\ref{fig:steps_tasks} maps the seven steps in the top row into five tasks in the bottom row: two AI chains (a length-three chain containing Steps 4--6 and a length-one chain containing only Step~2) and three manual steps.



\subsection{Short-Run Production}
\label{sec:shortrun.production}

The firm employs human labor to perform tasks and produce the output, and labor is the only input used in production.\footnote{
Viewed as a production function, our sequential technology is Leontief at the level of the individual firm: producing a unit of output requires completing every step of the workflow, in fixed proportions.
}
The worker carries out the entire step sequence $\mathcal{S} = (s_1, \dotsc, s_m)$, which the firm's AI strategy organizes into the task sequence $\T = (T_1, \dotsc, T_n)$, and is paid at the normalized rate of $w = 1$ per unit of time for the whole time production takes.
The cost of production is therefore the total expected time requirement of the workflow's tasks, and the firm chooses the AI deployment strategy that minimizes
\begin{align}
\label{eq:totalcost_shortterm}
\min_{\T} \ \sum_{T_b \in \T} \timecost{b}.
\end{align}

Problem \eqref{eq:totalcost_shortterm} is the firm's short-run cost minimization problem.
In the short run, AI is a purely productivity-improving tool, compressing the time a fixed set of steps takes by letting the firm optimize production cost through a combination of augmented and automated steps.

AI chaining makes solving this problem non-trivial.
Because a chain links several steps into a single unit of execution, the firm cannot solve \eqref{eq:totalcost_shortterm} step by step, and cost minimization instead requires comparing all feasible arrangements of the workflow.
This linkage gives rise to rich dynamics in how AI is deployed across production, and hence in how automation affects it.
We discuss these implications next.
